\documentclass[12pt,a4paper]{article}
\usepackage[utf8]{inputenc}
\usepackage[spanish,es-noshorthands]{babel}
\usepackage{amsmath}
\usepackage{siunitx}
\usepackage{circuitikz}
\usepackage[hidelinks]{hyperref}
\usepackage{geometry}
\geometry{a4paper, margin=2.5cm}

\title{Manual de Arquitectura y Estructura del Proyecto}
\author{Agente IA (Orquestador)}
\date{\today}

\begin{document}
\maketitle
\tableofcontents
\newpage

\section{Introducción}
Este manual describe la estructura de directorios y archivos presentes en el espacio de trabajo de Electrónica. El entorno está diseñado bajo una arquitectura de tres capas, que separa la lógica probabilística (LLMs) de la ejecución determinista (scripts), y también contiene recursos para IoT, Diseño de Circuitos Integrados (EDA) y Redacción Académica.

\section{Directorios Principales}
\begin{itemize}
    \item \textbf{\texttt{.agent/}}: Contiene las instrucciones del sistema para el agente de IA, definiendo su marco operativo, restricciones de entorno y protocolos de actuación.
    \item \textbf{\texttt{directives/}}: Capa 1 de la arquitectura. Contiene manuales de operación (SOPs) en formato YAML. Definen \textit{qué} debe hacerse (ej. \texttt{scrape\_website.yaml}, \texttt{git\_update.yaml}).
    \item \textbf{\texttt{execution/}}: Capa 3 de la arquitectura. Contiene scripts deterministas en Python con una sola responsabilidad (ej. \texttt{env\_diagnostic.py}, \texttt{alert\_user.py}).
    \item \textbf{\texttt{.tmp/}}: Directorio temporal utilizado por la capa de orquestación para almacenar estados de ejecución (como \texttt{run\_state.json}) y artefactos temporales de procesamiento.
    \item \textbf{\texttt{docs/}}: Documentación técnica y académica (19 subdirectorios: CIRC\_DISP\_ELECT/, SMPS/, EDA/, PROTECTOR\_120VAC/, ZBAR\_PRACTICAS/, EASYEDA/, RAG/, vLLM/, OPENCODE/, MEDIDOR\_ENERGIA/, PC\_ASUS/, PC\_PARA\_IA/, SERVIDOR\_POWEREDGE\_R610/, SISTEMA\_INTERNAC/, CIRC\_PARA\_RESP\_RAPIDAS/, PROYECTO\_MECATRONICA/, Ventilador\_3\_Velocidades/, entre otros).
    \item \textbf{\texttt{cursos/}}: Material didáctico, notas de cursos de Electrónica, apuntes de Laboratorio y Tesis (6 subdirectorios: DISP\_ELECTRONICOS/, INT\_ELECTRONICA/, TESIS/, PLAN\_ESTUDIOS/, LABORATORIO\_I\_FISICA/, LABORATORIO\_II\_FISICA/).
    \item \textbf{\texttt{chroma\_db/}}: Base de datos vectorial autogenerada, utilizada por el sistema RAG (Retrieval-Augmented Generation) para la memoria y búsqueda semántica del agente.
    \item \textbf{\texttt{Agente\_EDA/}}: Recursos específicos para el agente EDA, como esquemas (schemas) de validación y pruebas para la generación de netlists/BOM hacia EasyEDA.
\end{itemize}

\section{Archivos de Configuración y Scripts Raíz}
\begin{itemize}
    \item \textbf{\texttt{AGENTS.md}}: Reglas globales del proyecto, convenciones de código y estructura general del espacio de trabajo.
    \item \textbf{\texttt{.env}}: Plantilla de variables de entorno con credenciales comentadas (\texttt{GROQ\_API\_KEY}, \texttt{GITHUB\_TOKEN}) como referencia para la configuración local.
    \item \textbf{\texttt{.groq\_api\_key}}: Archivo oculto con la clave API de Groq (excluido de git vía \texttt{.gitignore}), usado por \texttt{rag\_system.py} y \texttt{agent\_eda.py}.
    \item \textbf{\texttt{.gitignore}}: Exclusiones de seguimiento git: \texttt{chroma\_db/}, entornos virtuales, \texttt{\_\_pycache\_\_}, archivos auxiliares LaTeX, clave API de Groq.
    \item \textbf{\texttt{requirements.txt}}: Dependencias para los scripts deterministas de ejecución: \texttt{psutil} y \texttt{PyYAML}.
    \item \textbf{\texttt{rag\_system.py}}: Script del Chatbot RAG. Vectoriza documentos (TEX, MD, PDF) en ChromaDB y gestiona las respuestas usando el LLM vía Groq.
    \item \textbf{\texttt{agent\_eda.py}}: Agente especializado en EDA. Extrae netlists y BOM de archivos LaTeX (circuitikz) para exportar hacia EasyEDA Standard.
    \item \textbf{\texttt{clean\_latex.py}}: Utilidad para eliminar archivos auxiliares y temporales de compilación de LaTeX (.aux, .log, .out, etc.).
    \item \textbf{\texttt{md\_to\_pdf.py} y \texttt{merge\_pdfs.py}}: Utilidades para la conversión y manipulación de documentos PDF.
    \item \textbf{\texttt{ren\_archivos.py}}: Script para el renombrado masivo de archivos eliminando cadenas específicas.
    \item \textbf{\texttt{test\_generator.py}}: Scripts de prueba (sin dependencias externas) para el generador JSON de EasyEDA.
    \item \textbf{\texttt{db\_state.json}}: Archivo que lleva el registro del estado y actualizaciones incrementales de la base de datos RAG vectorial.
    \item \textbf{\texttt{opencode.json}}: Archivo de configuración que enlaza las instrucciones del agente alojadas en \texttt{.agent/} (\texttt{\{"instructions": [".agent/*.md"]\}}).
    \item \textbf{\texttt{README.md}}: Documento de presentación del repositorio con instrucciones de instalación, uso y descripción general del proyecto.
    \item \textbf{\texttt{manual.tex}}: El presente manual en formato LaTeX, que describe la arquitectura, estructura y componentes del espacio de trabajo.
    \item \textbf{\texttt{git-update.sh} y \texttt{update\_repo.sh}}: Scripts de automatización en bash para sincronización y actualización del repositorio Git.
\end{itemize}

\section{Conclusión}
Esta estructura garantiza una alta modularidad. El agente IA actúa como "middleware" o capa de orquestación (Capa 2), leyendo intenciones, buscando las instrucciones en \texttt{directives/}, delegando el trabajo duro a \texttt{execution/} y documentando hallazgos en \texttt{chroma\_db/} o en archivos \texttt{.tex}. Todo ello cumpliendo con las reglas de codificación y formato estipuladas.

\end{document}
